\section{Introduction}
\label{sec:introduction}

\subsection{The Decadal Reapportionment Question}

The \dia{} specifies that the \ar{} algorithm is run fresh after each
decadal census, using the new \hh{} apportionment~\citep{huntington1928}
as the district count $k$ for each state.
This is the intended behavior: reapportionment is a constitutionally mandated
response to population change~\citep{wesberry1964}, and the algorithm
should reflect the new apportionment.

But what actually happens when a state gains or loses seats?
The \ar{} split tree is determined by the prime factorisation of $k$ together
with the prime fallback rule (T.4~\citep{b11paper}).
A change in $k$ can dramatically alter the tree topology---not because the
algorithm has changed, but because the arithmetic has.

This paper characterises an illustrative 2030 reapportionment shock.
We use a projection scenario~\citep{census2030proj} to identify states that
may gain or lose seats, trace how each seat change alters the \ar{} split
tree, and define the run package needed to measure tract-level boundary
movement.  The final 2030 apportionment will not be known until the official
enumeration and Huntington-Hill apportionment are complete.

\subsection{The Key Cases}

\paragraph{Texas ($38 \to 41$).}
$38 = 2 \times 19$ first splits into 19 two-district regions under the
largest-prime-first rule.
$41$ is prime, so the implemented prescription uses a binary $20:21$ fallback
at the root rather than a flat 41-way split.
This is the most dramatic factorisation change in the nation.
It is the priority case for package-backed boundary movement.

\paragraph{California ($52 \to 50$).}
$52 = 4 \times 13 = 2^2 \times 13$.
$50 = 2 \times 25 = 2 \times 5^2$.
Both are composite, but with different tree structures.
We characterise the structural change and expected boundary movement.

\paragraph{Florida ($28 \to 29$).}
$28 = 4 \times 7 = 2^2 \times 7$.
$29$ is prime: binary $14:15$ fallback, major structural change.

\paragraph{New York ($26 \to 25$).}
$26 = 2 \times 13$.
$25 = 5^2 = 5 \times 5$: balanced two-level tree.

\subsection{Main Finding}

The main finding is an evidence boundary: the reported Texas disruption
($d_{\rm Ham}=0.23$) is a plausible scenario value and a useful audit target,
but it must be archived with the same run metadata as other BISECT plan
packages before it can carry publication weight.  Algebraically, that value is
roughly nine one-step \recom{} units for $k=38$; this is a scale comparison,
not a measured mixing-time claim.

\subsection{Paper Structure}

Section~\ref{sec:projections} summarises 2030 reapportionment projections.
Section~\ref{sec:factorization} analyses the prime factorisation changes.
Section~\ref{sec:boundary} simulates boundary stability.
Section~\ref{sec:gerrychain-baseline} establishes the GerryChain baseline.
Section~\ref{sec:statutory} discusses the statutory implications.
Section~\ref{sec:conclusion} concludes.
